% =====================================================================
%  ANÁLISE DE PADRÕES E INCONSISTÊNCIA PARA IDENTIFICAÇÃO DE IMAGENS
%  GERADAS POR IA
%  Centro Universitário Fundação Santo André — 2026
%  Compilar com: pdflatex projeto-pesquisa.tex  (rodar 2x p/ o sumário)
% =====================================================================
\documentclass[12pt,a4paper,oneside]{article}

% ---------- Codificação e idioma ----------
\usepackage[utf8]{inputenc}
\usepackage[T1]{fontenc}
\usepackage[english,brazil]{babel}
\usepackage{lmodern}

% ---------- Layout ABNT (margens 3/2/3/2 cm) ----------
\usepackage[top=3cm,bottom=2cm,left=3cm,right=2cm]{geometry}
\usepackage{setspace}
\usepackage{indentfirst}
\usepackage{ragged2e}

% ---------- Recursos ----------
\usepackage{graphicx}
\usepackage{array}
\usepackage{longtable}
\usepackage{tabularx}
\usepackage{titlesec}
\usepackage{enumitem}
\usepackage{url}
\usepackage[hidelinks]{hyperref}

\graphicspath{{figuras/}}

% ---------- Seções não numeradas, em caixa alta, no sumário ----------
\newcommand{\secao}[1]{%
  \section*{\MakeUppercase{#1}}%
  \addcontentsline{toc}{section}{\MakeUppercase{#1}}%
}

\setlength{\parindent}{1.25cm}
\onehalfspacing

% ---------- Figura de anexo: \figuraanexo{nº}{legenda}{arquivo}{fonte} ----------
% Se o arquivo não existir em figuras/, imprime um espaço reservado.
\newcommand{\figuraanexo}[4]{%
  \begin{figure}[htb]
    \centering
    {\footnotesize Figura #1 -- #2\par}
    \vspace{2mm}
    \IfFileExists{figuras/#3}%
      {\includegraphics[width=0.75\textwidth]{#3}}%
      {\fbox{\parbox[c][4cm][c]{0.7\textwidth}{\centering\textit{[inserir arquivo figuras/#3]}}}}%
    \vspace{2mm}
    {\footnotesize\justifying Fonte: #4\par}
  \end{figure}%
}

\newcommand{\fontekingambit}{Gerado por Kingambit (2026), utilizando o
\textit{prompt}: ``Faça uma apresentação para estudantes do ensino médio com
base no projeto criado pelo grupo.''}

\begin{document}

% =====================================================================
%  ORDEM DO DOCUMENTO
%
%  Cada secao vive em seu proprio arquivo dentro de secoes/.
%  Assim duas pessoas editando partes diferentes nunca colidem no git.
%  \input insere o conteudo do arquivo aqui, como se estivesse digitado
%  neste ponto: a saida continua sendo um unico PDF.
%
%  Para compilar so a sua secao enquanto escreve (bem mais rapido),
%  comente as linhas \input das outras. Lembre de descomentar antes de
%  commitar.
% =====================================================================

% !TEX root = ../projeto-pesquisa.tex
% =====================================================================
%  CAPA
% =====================================================================
\begin{titlepage}
\begin{center}
\MakeUppercase{Centro Universitário Fundação Santo André}\\[2.5cm]

BIANCA GOMES DA SILVA\\
GIOVANNA DE MELLO DOROTEIO\\
IZABELA EIVE SILVA GASPARI\\ANALISE DE PADROES E INCONSISTENCIA PARA IDENTIFICAÇÃO DE IMAGENS GERADAS POR IA.
LEONARDO SIQUEIRA SOBRAL\\
LUIGI FAGE CESCHIM\\
LUISA PEREIRA DA SILVA ROMEIRO\\
MURILO AUGUSTO DA CRUZ\\
NATHAN GOMES DE PAULA\\
PEDRO RODRIGUES GODEC CARDOSO\\
RODRIGO FERREIRA FRANCISCO\\
VINICIUS HENRIQUE DOS SANTOS\\[3.5cm]

\textbf{\MakeUppercase{Análise de padrões e inconsistência para identificação
de imagens geradas por IA}}

\vfill
SANTO ANDRÉ\\
2026
\end{center}
\end{titlepage}

% =====================================================================
%  FOLHA DE ROSTO
% =====================================================================
\begin{titlepage}
\begin{center}
BIANCA GOMES DA SILVA\\
GIOVANNA DE MELLO DOROTEIO\\
IZABELA EIVE SILVA GASPARI\\
LEONARDO SIQUEIRA SOBRAL\\
LUIGI FAGE CESCHIM\\
LUISA PEREIRA DA SILVA ROMEIRO\\
MURILO AUGUSTO DA CRUZ\\
NATHAN GOMES DE PAULA\\
PEDRO RODRIGUES GODEC CARDOSO\\
RODRIGO FERREIRA FRANCISCO\\
VINICIUS HENRIQUE DOS SANTOS\\[3cm]

\textbf{\MakeUppercase{Análise de padrões e inconsistência para identificação
de imagens geradas por IA}}\\[3cm]
\end{center}

\hspace*{7cm}
\begin{minipage}{8cm}
\singlespacing\footnotesize\justifying
Projeto de Pesquisa apresentado como exigência parcial para a obtenção do grau
de bacharel, à Faculdade de Ciências Econômicas e Administrativas do Centro
Universitário Fundação Santo André.

\vspace{5mm}
Orientadora: Profa. Mestre Monica Walter Rodrigues.
\end{minipage}

\vfill
\begin{center}
SANTO ANDRÉ\\
2026
\end{center}
\end{titlepage}

% =====================================================================
%  FOLHA DE APROVAÇÃO
% =====================================================================
\begin{titlepage}
\singlespacing\justifying
\noindent\textbf{Candidatos:} Bianca Gomes da Silva, Giovanna de Mello Doroteio,
Izabela Eive Silva Gaspari, Leonardo Siqueira Sobral, Luigi Fage Ceschim,
Luisa Pereira da Silva Romeiro, Murilo Augusto da Cruz, Nathan Gomes de Paula,
Pedro Rodrigues Godec Cardoso, Rodrigo Ferreira Francisco, Vinicius Henrique
dos Santos.

\vspace{5mm}
\noindent\textbf{Título:} Análise de padrões e inconsistência para identificação
de imagens geradas por IA.

\vspace{5mm}
\noindent Projeto de Pesquisa apresentado como exigência parcial para a obtenção
do grau de bacharel, à Faculdade de Ciências Econômicas e Administrativas do
Centro Universitário Fundação Santo André, curso de Ciência de Dados e
Inteligência Artificial.

\vspace{10mm}
\noindent Data: \rule{1.5cm}{0.4pt}/\rule{1.5cm}{0.4pt}/\rule{2.5cm}{0.4pt}

\vspace{12mm}
\noindent Professor: \rule{6cm}{0.4pt}\hfill\rule{4cm}{0.4pt}

\noindent Instituição: \rule{6cm}{0.4pt}

\vspace{12mm}
\noindent Professor: \rule{6cm}{0.4pt}\hfill\rule{4cm}{0.4pt}

\noindent Instituição: \rule{6cm}{0.4pt}

\vspace{12mm}
\noindent Professor: \rule{6cm}{0.4pt}\hfill\rule{4cm}{0.4pt}

\noindent Instituição: \rule{6cm}{0.4pt}

\vspace{15mm}
\noindent Aprovado (\hspace{1cm})

\vspace{5mm}
\noindent Nota: \textbf{\rule{4cm}{0.4pt}}
\end{titlepage}


% !TEX root = ../projeto-pesquisa.tex
% =====================================================================
%  RESUMO
% =====================================================================
\secao{Resumo}
\justifying
O avanço das tecnologias e dos novos modelos generativos de inteligência
artificial tem apresentado um cenário desafiador no que tange à identificação de
imagens falsas, cuja sofisticação amplia o potencial de disseminação de
desinformação e golpes digitais no ambiente \textit{online}. Nesse contexto, o
presente projeto propõe o desenvolvimento de solução tecnológica voltada à
análise de imagens e capaz de apoiar a verificação da autenticidade e/ou
manipulação de arquivos digitais, tendo por objetivos: a) a implementação de
algoritmos de extração e análise de metadados e de padrões de compressão para
identificar assinaturas digitais características de processos de edição por IA;
b) a avaliação do potencial impacto do \textit{software} como mecanismo de
validação de conteúdo multimídia em larga escala na mitigação de desinformação;
c) a realização de testes comparativos de acurácia entre o \textit{software}
desenvolvido e o discernimento humano, estabelecendo métricas de desempenho para
a identificação de imagens sintéticas e d) a investigação das limitações da
técnica de \textit{Error Level Analysis} (ELA) frente ao refinamento de modelos
generativos, verificando a ocorrência de padronização de ruídos e compressão em
imagens de alta fidelidade. Para tanto, será realizada ação de extensão com cerca
de 40 estudantes de ensino médio de instituição de ensino em Santo André --- cuja
parceria ainda será formalizada, e que consistirá em atividade experimental com
os alunos, que divididos em dois grupos avaliarão cerca de 20 imagens
padronizadas (10 reais e 10 artificiais), sem e com o apoio da ferramenta
tecnológica desenvolvida para o projeto. Espera-se obter uma análise estatística
dos resultados com a identificação de padrões de erro, visando avaliar a eficácia
das técnicas computacionais frente à percepção humana, além de promover a
reflexão crítica sobre desinformação digital e o uso consciente de tecnologias.

\vspace{5mm}
\noindent\textbf{Palavras-chave:} análise de imagens, ciência de dados,
\textit{deepfakes}, desinformação, inteligência artificial.

% =====================================================================
%  ABSTRACT
% =====================================================================
\newpage
\secao{Abstract}
\justifying
\begin{otherlanguage}{english}
The advancement of technologies and new generative artificial intelligence models
has presented a challenging scenario regarding the identification of fake images,
whose sophistication increases the potential for the dissemination of
misinformation and digital scams in the online environment. In this context, this
project proposes the development of a technological solution focused on image
analysis and capable of supporting the verification of the authenticity and/or
manipulation of digital files, with the following objectives: a) the
implementation of metadata extraction and analysis algorithms and compression
patterns to identify digital signatures characteristic of AI editing processes;
b) the evaluation of the potential impact of the software as a mechanism for
large-scale multimedia content validation in mitigating misinformation; c) the
performance of comparative accuracy tests between the developed software and
human discernment, establishing performance metrics for the identification of
synthetic images; and d) the investigation of the limitations of the Error Level
Analysis (ELA) technique in the face of the refinement of generative models,
verifying the occurrence of noise standardization and compression in
high-fidelity images. To this end, an outreach activity will be carried out with
approximately 40 high school students from an educational institution in Santo
André --- whose partnership is yet to be formalized --- and will consist of an
experimental activity with the students, who, divided into two groups, will
evaluate approximately 20 standardized images (10 real and 10 artificial), both
with and without the support of the technological tool developed for the project.
The aim is to obtain a statistical analysis of the results, identifying error
patterns, in order to evaluate the effectiveness of computational techniques
compared to human perception, as well as promoting critical reflection on digital
misinformation and the conscious use of technology.

\vspace{5mm}
\noindent\textbf{Keywords:} artificial intelligence, data science, deepfakes,
disinformation, image analysis.
\end{otherlanguage}


% !TEX root = ../projeto-pesquisa.tex
% =====================================================================
%  LISTAS
% =====================================================================
\newpage
\secao{Lista de ilustrações}
\singlespacing
\begin{longtable}{@{}p{12cm}r@{}}
Figura 1 -- Imagem da resposta da parceria & 17\\
Figura 2 -- 1º Slide do programa & 18\\
Figura 3 -- 2º Slide do programa & 18\\
Figura 4 -- 3º Slide do programa & 19\\
Figura 5 -- 4º Slide do programa & 19\\
Figura 6 -- 5º Slide do programa & 20\\
Figura 7 -- 6º Slide do programa & 20\\
Figura 8 -- 7º Slide do programa & 21\\
Figura 9 -- 8º Slide do programa & 21\\
Figura 10 -- 9º Slide do programa & 22\\
Figura 11 -- 10º Slide do programa & 22\\
Figura 12 -- Abertura do Formulário de Pesquisa & 23\\
Figura 13 -- Imagem de cachorro gerada por IA & 24\\
Figura 14 -- Imagem de cachorro real & 25\\
Figura 15 -- Imagem de cachorro real & 26\\
Figura 16 -- Imagem de cachorro gerada por IA & 27\\
Figura 17 -- Imagem de macaco real & 28\\
Figura 18 -- Imagem de macaco gerada por IA & 29\\
Figura 19 -- Imagem de ornitorrinco gerada por IA & 30\\
Figura 20 -- Imagem de ornitorrinco real & 31\\
\end{longtable}

\vspace{6mm}
\secao{Lista de tabelas}
\begin{longtable}{@{}p{12cm}r@{}}
Cronograma de Execução & 15\\
\end{longtable}

\vspace{6mm}
\secao{Siglas}
\begin{longtable}{@{}p{2cm}p{11cm}@{}}
CNN  & Redes Neurais Convolucionais (\textit{Convolutional Neural Networks})\\
ELA  & Análise de Nível de Erro (\textit{Error Level Analysis})\\
EXIF & Formato de Arquivo de Imagem Intercambiável (\textit{Exchangeable Image
       File Format})\\
IA   & Inteligência Artificial\\
\end{longtable}
\onehalfspacing

% =====================================================================
%  SUMÁRIO
% =====================================================================
\newpage
\renewcommand{\contentsname}{Sumário}
\tableofcontents


% !TEX root = ../projeto-pesquisa.tex
% =====================================================================
%  INTRODUÇÃO
% =====================================================================
\newpage
\secao{Introdução}
\justifying
O avanço das tecnologias de inteligência artificial tem transformado
significativamente a produção e circulação de conteúdos digitais, especialmente
no que se refere à geração de imagens altamente realistas. Esse cenário amplia os
desafios relacionados à verificação da autenticidade de imagens digitais, uma vez
que modelos generativos são capazes de produzir conteúdos visualmente coerentes e
difíceis de distinguir de registros reais. Segundo Padilha \textit{et al.} (2021,
p. 131), o ambiente digital contemporâneo já é marcado por uma elevada presença
de conteúdos manipulados, incluindo montagens, informações fora de contexto e
materiais que dificultam a identificação de sua veracidade.

De acordo com as empresas DeepStrike (2025) e Bright Defense (2026), o volume de
\textit{deepfakes} (imagens, vídeos e áudios sintéticos) compartilhados
\textit{online} explodiu nos últimos anos: passou de aproximadamente 500 mil em
2023 para uma projeção de 8 milhões em 2025, representando um crescimento anual
de cerca de 1.500\%.

Nesse contexto, o presente projeto aborda a análise de padrões e inconsistências
na identificação de imagens geradas por inteligência artificial, com foco em
\textit{deepfakes} e demais formas de manipulação digital. A proposta consiste no
desenvolvimento de uma solução tecnológica voltada à análise de imagens, como uma
plataforma \textit{online}, capaz de apoiar a verificação da autenticidade de
arquivos digitais. Essa necessidade se intensifica diante da crescente
sofisticação dos modelos generativos, que ampliam o potencial de disseminação de
desinformação e golpes digitais no ambiente \textit{online}.

O problema central da pesquisa está relacionado à dificuldade de distinguir
imagens reais de imagens sintéticas, especialmente em casos de manipulação como o
\textit{copy-move}, em que a distribuição estatística dos \textit{pixels} pode
permanecer semelhante à original, dificultando a detecção visual e automatizada.
De acordo com Nagm \textit{et al.} (2024, p. 2), esse tipo de adulteração
representa um desafio relevante na área de forense digital, justamente por
preservar características estatísticas da imagem, o que compromete métodos
tradicionais de detecção baseados apenas em análise visual ou estatística
superficial.

Ainda nesse contexto, técnicas como a Análise de Nível de Erro (ELA) são
utilizadas para evidenciar inconsistências em imagens, por meio da amplificação
de diferenças de compressão que podem indicar possíveis manipulações (Nagm
\textit{et al.}, 2024, p. 4). No entanto, com o avanço dos modelos generativos de
alta fidelidade, a eficácia dessas abordagens tem sido reduzida, exigindo a
integração com métodos mais robustos, como redes neurais convolucionais (CNN),
capazes de identificar padrões complexos de manipulação, além da análise de
metadados EXIF, que fornecem informações técnicas sobre a origem e o histórico de
processamento das imagens (Yang \textit{et al.}, 2026, p. 15).

Estudos recentes mostram que a precisão humana na detecção de \textit{deepfakes}
fica em torno de 50--55\% (próxima ao nível de acaso), e apenas 0,1\% dos
participantes consegue identificar corretamente todos os conteúdos sintéticos e
reais apresentados (iProov, 2025).

Além disso, abordagens contemporâneas indicam que a autenticação de mídias
digitais não deve se basear apenas na análise do conteúdo visual, mas também na
verificação da proveniência e integridade dos arquivos. Segundo England
\textit{et al.} (2020, p. 115), em um cenário próximo, nem humanos nem algoritmos
serão capazes de distinguir de forma confiável conteúdos falsos de conteúdos
reais apenas pela aparência, tornando necessária a adoção de sistemas de
verificação baseados em origem, integridade e rastreabilidade da mídia.

A execução do projeto justifica-se pela necessidade de fortalecer a segurança
digital e promover o pensamento crítico no uso de mídias digitais. Dessa forma, o
projeto contribui para a formação acadêmica dos envolvidos, ao aplicar
conhecimentos de inteligência artificial, ciência de dados e segurança da
informação em um problema real, além de possuir relevância social ao auxiliar na
mitigação da desinformação e na promoção de um ambiente digital mais confiável.


% !TEX root = ../projeto-pesquisa.tex
% =====================================================================
%  TEMA
% =====================================================================
\secao{Tema}
Análise de padrões e inconsistência para identificação de imagens geradas por IA,
através de ELA e EXIF.

% =====================================================================
%  JUSTIFICATIVA
% =====================================================================
\secao{Justificativa}
O desenvolvimento deste projeto permite que os estudantes apliquem na prática
conhecimentos relacionados à tecnologia, análise de dados, segurança da
informação e inteligência artificial. Durante a pesquisa e construção do
aplicativo ou \textit{site}, os alunos poderão desenvolver habilidades como
investigação tecnológica, pensamento crítico, programação e análise de metadados
digitais.

Além disso, o projeto contribui para a formação acadêmica ao aproximar o
conhecimento teórico da aplicação prática voltada para problemas reais da
sociedade, como a disseminação de \textit{fake news} e conteúdos manipulados pela
inteligência artificial.

Para a faculdade, o projeto fortalece as atividades de extensão universitária,
promovendo a integração entre ensino, pesquisa e sociedade. A iniciativa também
incentiva o desenvolvimento de soluções tecnológicas inovadoras que podem gerar
impacto social positivo, além de estimular a produção de conhecimento e a
conscientização digital da comunidade.


% !TEX root = ../projeto-pesquisa.tex
% =====================================================================
%  OBJETIVOS
% =====================================================================
\secao{Objetivos}
\noindent\textbf{Objetivo Geral:} Verificar a viabilidade prática da
identificação de assinaturas digitais, metadados e padrões de compressão em
imagens, visando detectar conteúdos gerados por inteligência artificial e avaliar
seu impacto no contexto da desinformação digital.

\vspace{4mm}
\noindent\textbf{Objetivos Específicos:}
\begin{enumerate}[label=\arabic*~--,leftmargin=1.5cm,itemsep=4pt]
  \item Implementar algoritmos de extração e análise de metadados e padrões de
        compressão para identificar assinaturas digitais características de
        processos de edição por IA.
  \item Avaliar o impacto de ferramentas de verificação na mitigação de
        desinformação, analisando o potencial do \textit{software} como mecanismo
        de validação de conteúdos multimídia em larga escala.
  \item Realizar testes comparativos de acurácia entre o \textit{software}
        desenvolvido e o discernimento humano, estabelecendo métricas de
        desempenho para a identificação de imagens sintéticas.
  \item Investigar as limitações da técnica de \textit{Error Level Analysis}
        (ELA) frente ao refinamento de modelos generativos, verificando a
        ocorrência de padronização de ruídos e compressão em imagens de alta
        fidelidade.
\end{enumerate}


% !TEX root = ../projeto-pesquisa.tex
% =====================================================================
%  METODOLOGIA
% =====================================================================
\secao{Metodologia}
A metodologia consiste em uma pesquisa de abordagem quantitativa e experimental,
estruturada como um estudo de campo controlado, com o objetivo de avaliar a
eficácia de técnicas computacionais na identificação de imagens geradas por
inteligência artificial, especialmente \textit{deepfakes}. A investigação
baseia-se na comparação entre a percepção visual humana e o uso de ferramentas de
ciência de dados, como leitura de metadados e análise de padrões de compressão.

A ação de extensão será realizada com estudantes do ensino médio de instituições
da região de Santo André, envolvendo aproximadamente 40 participantes. A proposta
adota uma abordagem dialógica, na qual os alunos atuarão de forma ativa na
construção do conhecimento, compartilhando percepções, hipóteses e dificuldades
ao longo das atividades, promovendo a troca de saberes entre a equipe do projeto
e a comunidade escolar.

O projeto será desenvolvido em etapas interdependentes. Inicialmente, será feita a
curadoria de 20 imagens digitais de alta resolução, sendo 10 autênticas e 10
geradas por inteligência artificial. A seleção considerará critérios que aumentem
a complexidade da avaliação, como texturas, iluminação, profundidade de campo e
ausência de artefatos evidentes. Também serão padronizados formato do arquivo,
resolução e nível de compressão, garantindo consistência metodológica no
experimento.

Paralelamente, será desenvolvido um protótipo de ferramenta digital, na forma de
um \textit{site} ou aplicação simples, com funcionalidades voltadas à verificação
de imagens. Entre elas, destacam-se a extração de metadados EXIF, a identificação
de \textit{softwares} de edição, a análise de padrões de compressão e a aplicação
da técnica de Análise de Nível de Erro (ELA), permitindo evidenciar
inconsistências que podem indicar manipulação digital. O desenvolvimento
priorizará a clareza dos resultados, de modo que os participantes compreendam os
indícios apresentados.

Na sequência, será realizada a formalização ética, com autorização da instituição
de ensino e aplicação do Termo de Consentimento de Participação, assegurando
participação voluntária e tratamento anônimo dos dados para fins acadêmicos.

A aplicação ocorrerá em ambiente escolar, iniciando-se com uma breve introdução
sobre inteligência artificial, \textit{deepfakes} e desinformação digital, em
linguagem acessível, sem interferir no desempenho dos participantes. Em seguida,
os estudantes serão divididos aleatoriamente em dois grupos iguais: o primeiro
grupo realizará a avaliação apenas com base na percepção visual, enquanto o outro
grupo contará com apoio da ferramenta desenvolvida. Cada participante
classificará as 20 imagens como reais ou artificiais dentro de um tempo definido,
com respostas registradas em formulário digital, garantindo padronização das
condições da atividade.

Após a atividade, será realizada uma discussão coletiva, na qual serão
apresentados exemplos e resultados preliminares. Os estudantes poderão refletir
sobre seus acertos e erros, compartilhar estratégias e discutir dificuldades,
estimulando o pensamento crítico sobre a confiabilidade de conteúdos digitais e
os limites da percepção humana.

Posteriormente, os dados serão submetidos à análise estatística, comparando as
taxas de acerto entre os grupos e identificando padrões de erro e níveis de
dificuldade das imagens. Isso permitirá avaliar a eficácia das técnicas
computacionais e suas limitações frente ao avanço dos modelos generativos. Por
fim, os resultados serão organizados em relatórios e produções acadêmicas, com
gráficos e análises interpretativas. Também será realizada uma devolutiva à
instituição participante, com orientações sobre identificação de conteúdos
manipulados, reforçando o caráter educativo da proposta.


% !TEX root = ../projeto-pesquisa.tex
% =====================================================================
%  RECURSOS
% =====================================================================
\secao{Recursos}
Os recursos utilizados para a realização do projeto envolvem: \textit{notebooks}
de nova geração, conexão estável à rede de internet e a participação de 11
indivíduos para a construção do projeto e sua documentação.

A relação de parceria foi estabelecida no Colégio Pentágono, localizado na cidade
de Santo André, São Paulo, ambiente este adequado para a aplicação prática do
projeto.

% =====================================================================
%  CRONOGRAMA
% =====================================================================
\secao{Cronograma de execução}
\singlespacing
\begin{center}
\begin{tabular}{|>{\raggedright\arraybackslash}p{5.5cm}|c|c|c|c|c|}
\hline
\textbf{AÇÃO} & \textbf{Ago.} & \textbf{Set.} & \textbf{Out.} & \textbf{Nov.} &
\textbf{Dez.}\\
\hline
Desenvolvimento da aplicação do projeto & & & & & \\
\hline
Ajustes e melhorias do projeto (\textit{app}) & & & & & \\
\hline
Apresentação do projeto em escola parceira & & & & & \\
\hline
Validação dos resultados obtidos & & & & & \\
\hline
Execução e bateria de testes do projeto & & & & & \\
\hline
Análise dos resultados das baterias de testes & & & & & \\
\hline
Montagem da apresentação de ACEx & & & & & \\
\hline
Apresentação do projeto final (SAPex) & & & & & \\
\hline
\end{tabular}
\end{center}
\onehalfspacing


% !TEX root = ../projeto-pesquisa.tex
% =====================================================================
%  REFERÊNCIAS
% =====================================================================
\newpage
\secao{Referências bibliográficas}
\singlespacing
\begin{flushleft}
BRIGHT DEFENSE. \textbf{150+ Deepfake statistics (March 2026)}. 2026. Disponível
em: \url{https://www.brightdefense.com/resources/deepfake-statistics/}. Acesso
em: 12 abr. 2026.

\vspace{4mm}
DEEPSTRIKE. \textbf{Deepfake statistics 2025: AI Fraud Data \& Trends}. 2026.
Disponível em: \url{https://deepstrike.io/blog/deepfake-statistics-2025}. Acesso
em: 12 abr. 2026.

\vspace{4mm}
ENGLAND, Paul \textit{et al}. AMP: Authentication of media via provenance. In:
\textbf{\textit{Proceedings of the 12th ACM Multimedia Systems Conference}}.
2021. p. 108-121. Disponível em: arXiv. Acesso em: 18 mar. 2026.

\vspace{4mm}
IPROOV. \textbf{Study reveals deepfake blindspot}: most consumers can't identify
AI-generated content. [S.l.], 2024. Disponível em:
\url{https://www.iproov.com/press/study-reveals-deepfake-blindspot-detect-ai-generated-content}.
Acesso em: 29 maio 2026.

\vspace{4mm}
NAGM, Ahmad M. \textit{et al}. Detecting image manipulation with ELA-CNN
integration: a powerful framework for authenticity verification.
\textbf{\textit{PeerJ Computer Science}}, v. 10, p. e2205, 2024. Disponível em:
\url{https://peerj.com/articles/cs-2205/}. Acesso em: 26 mar. 2026.

\vspace{4mm}
PADILHA, Rafael \textit{et al}. A Inteligência Artificial e os desafios da Ciência
Forense Digital no século XXI. \textbf{Estudos Avançados}, v. 35, n. 101,
p. 113-138, 2021. Disponível em: SciELO. Acesso em: 18 mar. 2026.

\vspace{4mm}
YANG, Pengpeng \textit{et al}. Forensic Analysis for Source Camera Identification
from EXIF Metadata. \textbf{\textit{Journal of Imaging}}, v. 12, n. 3, p. 110,
2026. Disponível em: MDPI. Acesso em: 18 mar. 2026.
\end{flushleft}
\onehalfspacing


% !TEX root = ../projeto-pesquisa.tex
% =====================================================================
%  ANEXOS
% =====================================================================
\newpage
\secao{Anexo I --- Confirmação da parceria}
\figuraanexo{1}{Imagem da resposta da parceria}{fig01-parceria.png}%
{Conversa realizada por e-mail entre o integrante do projeto Nathan Gomes Paula e
o professor da instituição de ensino Pentágono, Rafael Cirera.}

\newpage
\secao{Anexo II --- Carta convite enviada}
% [Inserir aqui o conteúdo da carta convite]

\newpage
\secao{Anexo III --- Modelo de apresentação para pesquisa}

\figuraanexo{2}{1º Slide do programa}{fig02-slide01.jpg}{\fontekingambit}
\figuraanexo{3}{2º Slide do programa}{fig03-slide02.jpg}{\fontekingambit}
\figuraanexo{4}{3º Slide do programa}{fig04-slide03.jpg}{\fontekingambit}
\figuraanexo{5}{4º Slide do programa}{fig05-slide04.jpg}{\fontekingambit}
\figuraanexo{6}{5º Slide do programa}{fig06-slide05.jpg}{\fontekingambit}
\figuraanexo{7}{6º Slide do programa}{fig07-slide06.jpg}{\fontekingambit}
\figuraanexo{8}{7º Slide do programa}{fig08-slide07.jpg}{\fontekingambit}
\figuraanexo{9}{8º Slide do programa}{fig09-slide08.jpg}{\fontekingambit}
\figuraanexo{10}{9º Slide do programa}{fig10-slide09.jpg}{\fontekingambit}
\figuraanexo{11}{10º Slide do programa}{fig11-slide10.jpg}{\fontekingambit}

\newpage
\secao{Anexo IV --- Modelo de questionário com imagens para pesquisa}

\figuraanexo{12}{Abertura do Formulário de Pesquisa}{fig12-formulario.png}%
{Elaborado pelo grupo via Google Forms.}

\figuraanexo{13}{Imagem de cachorro gerada por IA}{fig13-cachorro-ia.jpg}%
{Gerado por Google Gemini, utilizando o \textit{prompt}: ``Foto realista de um
Westie branco peludo em um jardim, de pé em meio a flores silvestres e grama
alta. A luz solar dourada realça o pelo e a cauda levantada. O cachorro está
usando uma \textit{tag} `Oliver'. \textit{Bokeh} suave ao fundo''.}

\figuraanexo{14}{Imagem de cachorro real}{fig14-cachorro-real.jpg}%
{GaúchaZH (2025).}

\figuraanexo{15}{Imagem de cachorro real}{fig15-cachorro-real.jpg}%
{Canil Sonhare (s.d.).}

\figuraanexo{16}{Imagem de cachorro gerada por IA}{fig16-cachorro-ia.jpg}%
{Gerado por Google Gemini, utilizando o \textit{prompt}: ``Um cachorrinho de
\textit{Dachshund dapple} fofo sentado na varanda de uma casa, com sua guia e
\textit{tag} de identificação penduradas. Foto íntima focada na expressão
inocente do filhote. O pelo mesclado e as orelhas longas são proeminentes''.}

\figuraanexo{17}{Imagem de macaco real}{fig17-macaco-real.jpg}%
{La Crónica de Hoy (2025).}

\figuraanexo{18}{Imagem de macaco gerada por IA}{fig18-macaco-ia.jpg}%
{Gerado por Google Gemini, utilizando o \textit{prompt}: ``\textit{Close-up}
médio de um macaco-prego capuchinho em seu habitat natural, equilibrado em um
galho de árvore. O pelo marrom e o rosto escuro são nítidos contra o verde
vibrante da selva desfocada. Iluminação ambiente e natural''.}

\figuraanexo{19}{Imagem de ornitorrinco gerada por IA}%
{fig19-ornitorrinco-ia.jpg}%
{Gerado por Google Gemini, utilizando o \textit{prompt}: ``\textit{Close-up}
detalhado do rosto e do bico de um ornitorrinco isolado em branco. A textura do
bico e do pelo. Posição de pé, olhando para a frente. Imagem macro, nítida e bem
iluminada''.}

\figuraanexo{20}{Imagem de ornitorrinco real}{fig20-ornitorrinco-real.jpg}%
{MHI Omics (s.d.).}



\end{document}
